% =========================================================
% File: frontmatter/abstract.tex
% =========================================================
\chapter*{Abstract}
\addcontentsline{toc}{chapter}{Abstract}

This study addresses the problem of vision-based running technique classification utilizing unconstrained, in-the-wild internet videos. Traditional biomechanical analysis relies heavily on marker-based motion capture, which is restricted to controlled laboratory settings. To bridge this research gap, we propose an automated computer vision pipeline that processes raw internet running videos using markerless pose estimation via MediaPipe. The structural framework extracts a custom spatial-temporal signal \code{contact_diff} through sequential landmark filtering and normalization. An autocorrelation-based automatic gait cycle detection system is established to isolate distinct strides, which are subsequently downsampled into a standardized 8-frame gait phase representation.

The classification task evaluates and compares landmark-based classifiers alongside an image-based MobileNetV2 + LSTM pipeline. To diagnose performance bottlenecks, comprehensive model, signal, and feature ablation studies are performed against a handmade cycle segmentation upper bound and an automatic baseline. While the system demonstrates that markerless pose sequences can effectively encode distinct running-technique characteristics when cycles are isolated cleanly, comprehensive error analysis reveals significant performance degradation in fully automated end-to-end extraction. Specifically, the handmade upper-bound pipeline achieves a test macro-F1 of 0.8991 and an accuracy of 0.9107 on the controlled test split, whereas the fully automatic autocorrelation-based baseline reaches a macro-F1 of 0.5417 and an accuracy of 0.5455. A further 5-fold video-level cross-validation robustness check confirms the same trend under a shared MLP protocol: handmade cycles achieve 0.860 $\pm$ 0.016 macro-F1, while autocorr cycles achieve 0.615 $\pm$ 0.081, leaving a fair gap of 0.245 macro-F1. A 100-permutation test on the autocorr MLP setting does not reach statistical significance ($p = 0.2475$), suggesting that the autocorr-derived features remain noisy and are not yet robustly separable under a simple classifier. The dataset contains 75 videos and 152 extracted running cycles, but the upper-bound and automatic baseline metrics are reported on their respective controlled test splits rather than on the full dataset. The empirical findings indicate that future deployment readiness depends fundamentally on improving upstream cycle-boundary precision rather than classifier architectural complexity.

\vspace{1.0cm}
\noindent \textbf{Keywords:} Markerless Pose Estimation, Action Classification, Gait Cycle Detection, Deep Learning, Biomechanics.
